\documentclass{resume}

\usepackage{zh_CN-Adobefonts_external}
\usepackage{enumitem}
\usepackage{linespacing_fix}
\usepackage{cite}

\begin{document}
\pagenumbering{gobble}

\noindent
\name{李昱辰}
\otherInfo{27 届}{意向岗位：Java 后端开发实习 / Agent 开发实习}{}{}
\otherInfo{可尽快到岗 \textperiodcentered\ 每周可实习 5 天 \textperiodcentered\ 可连续实习 6 个月}{邮箱：edllyc041202@gmail.com}{}{}

\vspace{-0.8em}

\section{教育背景}
\datedsubsection{\textbf{安徽大学} - 数字媒体技术专业 - 本科}{2023.09 - 2027.06}
\begin{itemize}[parsep=0.5ex]
    \item \textbf{科研与竞赛成果}：以学生一作身份向 \textit{IEEE TMI} 在投算法论文 1 篇（外审），提出基于拓扑学先验的连通性损失函数；获美赛二等奖、数学建模省赛三等奖、三等学业奖学金。
    \item \textbf{外语能力}：雅思（IELTS）7.0 分，具备英文技术文档、学术论文阅读与复现能力。
    \item \textbf{成绩情况}：均分 87.6/100，系统修读数据结构、操作系统、计算机网络、数据库等计算机核心课程。
\end{itemize}

\section{专业技能}
\begin{itemize}[parsep=0.5ex]
    \item \textbf{Java / JVM / 并发}：掌握 Java 核心基础，熟悉集合、并发编程与 JVM 内存模型，具备基于 CompletableFuture 的异步编排与高并发链路设计经验。
    \item \textbf{Spring Boot / MyBatis-Plus / MySQL}：熟练使用 Spring Boot、MyBatis-Plus、MySQL 进行后端服务开发，具备索引优化、事务隔离、SQL 调优与业务链路落库实践经验。
    \item \textbf{Redis / Redisson / Lua / RocketMQ / Sentinel / Docker / Maven / Git}：掌握 Redis/Redisson/Lua 缓存与并发控制方案，具备多级缓存、分布式锁、库存原子扣减、MQ 事务消息/延迟消息、限流熔断治理经验；熟悉 Docker 基础容器化部署以及 Maven、Git 日常工程协作。
    \item \textbf{LangChain4j / WebFlux / Reactor / RAG / Ollama / RAGAS / LLM-as-a-Judge}：具备 LangChain4j + Spring WebFlux + Project Reactor 智能体应用开发经验，掌握 AiServices 驱动的实体提取、查询改写、意图路由与复杂度判定；熟悉 Plan-and-Execute、Parent-Child、Milvus/OpenSearch 混合检索、RRF、BM25+、Reranker、Planned RAG，并可基于 Ollama、RAGAS、LLM-as-a-Judge 完成本地辅助判定与离线评测。
\end{itemize}

\section{项目经历}

\datedsubsection{\textbf{12306 高并发售票核心网关}（后端架构开发）}{2026.01 - 2026.05}
\textbf{项目技术}：Spring Boot, Redis, Redisson, Caffeine, RBloomFilter, RocketMQ, Sentinel, MyBatis-Plus\\
\textbf{项目描述}：面向 12306 高并发售票场景，围绕查票、订票、退款等核心链路完成缓存治理、库存预扣、事务一致性与补偿闭环设计。\\
\textbf{项目职责}：
\begin{enumerate}[parsep=0.5ex]
    \item \textbf{缓存与高并发治理}：基于 Caffeine + Redis + RBloomFilter 构建多级缓存，结合 Redisson 分布式锁与基于 CompletableFuture 的 singleflight 回源合并治理热点 Key；在查票热路径压测中，800 并发下达 5409.54 QPS、0\% HTTP 错误率，P95 344ms。
    \item \textbf{库存扣减与分桶模型}：基于 Redis Lua 实现库存原子预扣减，设计 Segmented Stock 分桶库存模型分散热点流量，降低高并发订票场景下的库存争抢与超卖风险。
    \item \textbf{订单异步化与事务消息}：通过 RocketMQ 事务消息 + TransactionListener 串联 Redis 预扣、订单落库与 MySQL 扣减，构建异步削峰与最终一致性链路；在 200 用户订票压测中，达 510.94 TPS、0\% 错误率，P95 318ms，未观察到超卖且异步收敛后一致性校验通过。
    \item \textbf{退款补偿与稳定性保护}：围绕退款链路设计 refund 字段与 Redis Lua 幂等回补机制，解决 MQ 重试导致的重复补库存问题；基于 RocketMQ 延迟消息实现超时未支付订单自动关单，并结合 Sentinel 完成限流与热点保护。
\end{enumerate}

\datedsubsection{\textbf{12306 智能数字员工}（Reactive Multi-Agent / RAG 智能体研发）}{2026.01 - 2026.05}
\textbf{项目技术}：Spring WebFlux, Project Reactor, LangChain4j, Milvus, OpenSearch, 智谱 GLM, Ollama, RAGAS\\
\textbf{项目描述}：面向 12306 规章问答与购退改查场景，构建响应式多智能体与 RAG 执行链路，覆盖意图路由、规划执行、混合检索、记忆管理与流式反馈。\\
\textbf{项目职责}：
\begin{enumerate}[parsep=0.5ex]
    \item \textbf{多专家意图路由}：设计多专家意图打分机制，对 ACTION / TICKET / RAG / CHITCHAT 四类意图并行评分，结合最高分、分差与置信度完成分层路由决策，降低动作请求、票务查询与政策问答之间的误路由。
    \item \textbf{政策路由与语义桥接}：基于 AiServices 构建上下文依赖判定、政策分类与政策证据判定链路，结合 Session 前缀隔离、ConversationBridge 与本地部署的 Ollama 轻模型辅助判定，提升追问承接与跨专家切换时的语义连续性。
    \item \textbf{Plan-and-Execute 执行链路}：将传统 ReAct 升级为 Plan-and-Execute，由 Planner 输出 JSON 计划拆解复杂购退票请求，Executor 调用票务工具并消费 ToolObservation 结构化观测结果，结合 Replan 机制处理执行异常。
    \item \textbf{混合检索与 Planned RAG}：采用 Parent-Child 切片与 Milvus Dense + OpenSearch BM25 双路召回，结合 RRF 融合、BM25+ 打分与 Reranker 精排优化复杂政策问答命中质量，并对复杂问题引入子问题拆解实现 Planned RAG。
    \item \textbf{评测、流式交互与记忆}：基于评测集执行 RAGAS 评测，结果达到 context precision 0.881、context recall 0.935；同时基于 WebFlux + Project Reactor + Sinks.Many 实现流式输出，并接入 Redis 分布式语义缓存与短期滑动窗口 + 长期摘要压缩双层记忆。
\end{enumerate}

\end{document}
